\subsection{Náročnost provozu}
\label{subsec:criteria-operational}

Náročnost provozu určuje, jaké zdroje jsou potřeba k dlouhodobému provozování modelu v produkčním prostředí.
Zahrnuje jak finanční náklady, tak nároky na infrastrukturu a údržbu.

U cloudových API jsou náklady typicky určeny cenovými modely a jak bylo popsáno v \secref[Sekci]{subsec:llm-tokens-parameters-context}, jsou založeny na počtu zpracovaných tokenů modelem.
Pro typickou analýzu jednoho studentského odevzdání, zahrnující systémový prompt, zadání a zdrojový kód, je potřeba přibližně 2~000--4~000 vstupních tokenů a 500--2~000 výstupních tokenů při použití one-shot (OS) logiky.
Pro použití chain-of-thought (COT) přístupu, který zahrnuje více kroků a tedy více požadavků na model, se tyto počty mohou zvýšit až na 8~000--16~000 vstupních tokenů a 4~000--8~000 výstupních tokenů.
Celkové náklady na jedno odevzdání tak závisí na zvoleném modelu a jeho cenové politice.

Pro on-premise nasazení jsou náklady primárně jednorázové a zahrnují pořízení hardware, především tedy grafických karet s dostatečnou kapacitou VRAM pro efektivní inference, a náklady na údržbu a aktualizace modelu které jsou průběžně mnohem nižší než u cloudových API, ale vyžadují technickou odbornost pro správu a optimalizaci.

\endinput